\section{Tesla Optimus: observaciones desde la investigación y el desarrollo de primera línea}

\subsection{Experiencia de trabajo con Elon Musk}

Yang Shuo (杨硕) se incorporó al equipo de Optimus de Tesla después de doctorarse en 2023 y trabajó ahí alrededor de un año y medio. Sobre el estilo de trabajo de Elon Musk, hizo las siguientes observaciones:

\begin{itemize}
    \item \textbf{Concentración extrema}: durante las juntas prácticamente no se distrae ni pierde el hilo; observa con mucha atención a quien presenta y sigue todo el proceso
    \item \textbf{Eficiencia extrema para recabar información}: recoge los avances técnicos directamente de los ingenieros de primera línea; cada ingeniero pasa a exponer unos minutos y luego sigue el siguiente, de modo que en hora y media puede escuchar los reportes de veinte o treinta ingenieros
    \item \textbf{Decisiones a nivel macro}: al final, desde una perspectiva macro, asigna a los responsables del proyecto los ajustes de dirección general, en lugar de enfrascarse en los detalles técnicos
    \item \textbf{Paciencia con el proyecto Optimus}: sabe que cada detalle está en la frontera del conocimiento y que la certeza es baja, y tiene un buen control de conjunto sobre todo el proceso
\end{itemize}

Sobre el rumor de que «despiden a la gente durante los reportes», Yang Shuo señaló que esto no ha ocurrido en el proyecto Optimus; entiende que podría suceder en proyectos que atraviesan un momento crítico o que requieren decisiones de gran envergadura.

\begin{knowledgebox}{La ventaja sistémica de Tesla Optimus}
Yang Shuo considera que Tesla es, en cuanto a confiabilidad de robots, lo más cercano al nivel de Boeing; la razón central es que su negocio de robótica heredó la infraestructura completa ya establecida por la línea de negocio automotriz. Si Tesla fabrica 5,000 robots en un año, sin importar dónde se desplieguen, ya cuenta con un sistema y un proceso muy maduros para recolectar y analizar fallas. Esta capacidad sistémica es algo que una startup común difícilmente puede replicar en el corto plazo.
\end{knowledgebox}

\subsection{El sentido de misión en DJI y Tesla}

Yang Shuo sintió profundamente en DJI y en Tesla una cultura empresarial impulsada por un \textbf{sentido de misión}. Los puntos en común de ambas empresas son:

\begin{itemize}
    \item Buscan constantemente crear categorías de producto y experiencias de usuario completamente nuevas
    \item DJI convirtió los drones y las cámaras portátiles en categorías totalmente nuevas
    \item Tesla perfeccionó y renovó el automóvil eléctrico; internamente, Elon insiste una y otra vez en el sentido de misión de «cambiar la forma en que la humanidad se desplaza en el futuro»
    \item Al desarrollar robots humanoides, Elon siempre piensa el problema desde la perspectiva del «futuro de toda la humanidad»
\end{itemize}

Yang Shuo considera que este sentido de misión es una razón importante del éxito de ambas empresas: aunque el valor comercial se manifiesta después, es precisamente el sentido de misión lo que impulsó el nacimiento de estos productos innovadores.

\subsection{Avances técnicos de Optimus}

Dentro de Tesla, Yang Shuo ya ha visto indicios de que el robot cuenta con varias capacidades básicas, entre ellas: un desempeño básico, caminar de forma simple, doblar ropa por teleoperación, y recoger y colocar objetos de manera automática. Estima que en menos de cinco años, o incluso en uno o dos años, estas capacidades podrían convertirse en productos con valor de aplicación real.

Pero también señala con claridad que aún hay una brecha considerable entre «lograr hacerlo» y «poder venderlo», y que la clave está en el escenario de uso. Por ejemplo, tomar una botella de bebida y tomar una pieza pequeña de Lego son tareas de dificultad completamente distinta; si el escenario de aplicación es tomar y colocar objetos relativamente grandes (como una caja de paquetería de Amazon), el robot de Agility Robotics ya ha demostrado seis o siete horas continuas moviendo cajas sin ningún fallo, lo cual está muy cerca de la aplicación comercial. Pero si lo que se espera es «comprar un robot, ponerlo en casa y que con solo decir una frase recoja la mesa», esto probablemente no se logre dentro de cinco a diez años.

Yang Shuo también señaló una observación interesante: todos los especialistas en robótica, al predecir el futuro, dicen «cinco años», pero el desarrollo real suele tardar más. Sin embargo, los avances que él mismo ha visto dentro de Tesla lo llevan a creer que, en ciertos escenarios específicos, los productos robóticos realmente podrían necesitar solo uno o dos años.

\begin{knowledgebox}{La inercia psicológica de la «predicción a cinco años»}
Yang Shuo señala un fenómeno interesante de la industria: a casi todos los especialistas en robótica, cuando se les pregunta cuándo los robots entrarán a los hogares, responden «cinco años». Esto refleja una inercia psicológica: cinco años es un plazo lo bastante lejano como para no tener que responder por él, pero lo bastante cercano como para generar entusiasmo. Aunque Yang Shuo también dio un marco temporal similar, basándose en sus observaciones reales dentro de Tesla mantiene una postura más optimista respecto a escenarios verticales específicos, y a la vez conserva cautela respecto al escenario de un asistente doméstico de propósito general.
\end{knowledgebox}

\subsection{Resumen de la sección}

La experiencia en Tesla le permitió a Yang Shuo ser testigo de primera mano de cómo un sistema industrial a gran escala puede potenciar la investigación y el desarrollo robótico de vanguardia. El estilo eficiente de Elon para recabar información y tomar decisiones a nivel macro, la infraestructura sistémica que Tesla heredó de su negocio automotriz, y la cultura empresarial impulsada por el sentido de misión, son todas fuentes de la ventaja competitiva particular del proyecto Optimus.
